
\documentclass{article}
\usepackage{colortbl}
\usepackage{makecell}
\usepackage{multirow}
\usepackage{supertabular}

\begin{document}

\newcounter{utterance}

\centering \large Interaction Transcript for game `ta\_mastermind', experiment `Mastermind{-}v0', episode 9 with BSU{-}SLIM\_\_Qwen3.5{-}9B{-}Base\_\_playpen{-}prm{-}9b{-}best\_of\_n.
\vspace{24pt}

{ \footnotesize  \setcounter{utterance}{1}
\setlength{\tabcolsep}{0pt}
\begin{supertabular}{c@{$\;$}|p{.15\linewidth}@{}p{.15\linewidth}p{.15\linewidth}p{.15\linewidth}p{.15\linewidth}p{.15\linewidth}}
    \# & \multicolumn{2}{c}{Player} && \multicolumn{2}{c}{Game Master} \\
    \hline

    \theutterance \stepcounter{utterance}  
    & & & \multicolumn{4}{p{0.6\linewidth}}{
        \cellcolor[rgb]{0.9,0.9,0.9}{
            \makecell[{{p{\linewidth}}}]{
                \texttt{\tiny{[P1$\langle$GM]}}
                \texttt{You are a competitive game player. Make sure you read the game instructions carefully, and always follow the required format.} \\
\\ 
\texttt{{[}GAME{]} You are playing Mastermind.} \\
\texttt{You need to find the code that is 4 digits long, each digit from 1 to 6, with no duplicates.} \\
\texttt{In your response, you can submit your guess in the following format: '{[}2 1 4 5{]}'.} \\
\texttt{After each guess, you will receive feedback in the form of black and white pegs.} \\
\texttt{A black peg indicates a correct digit in the correct position, while a white peg indicates a correct digit in the wrong position.} \\
\texttt{You have 20 turns to guess the code.} \\
\\ 
            }
        }
    }
     \\ \\

    \theutterance \stepcounter{utterance}  
    & \multicolumn{4}{p{0.6\linewidth}}{
        \cellcolor[rgb]{0.561,0.737,0.561}{
            \makecell[{{p{\linewidth}}}]{
                \texttt{\tiny{[P1$\rangle$GM]}}
                \texttt{I'm ready to begin! Please provide the initial feedback (black and white pegs) for my first guess, or let me know if you'd like me to submit a random valid guess to start the game.} \\
\\ 
\texttt{Remember:} \\
\texttt{{-} The code is 4 digits long.} \\
\texttt{{-} Each digit is from 1 to 6.} \\
\texttt{{-} No duplicates are allowed.} \\
\texttt{{-} I'll submit guesses in the format: '{[}2 1 4 5{]}'.} \\
\texttt{{-} Feedback will be given as black and white pegs.} \\
\texttt{{-} I have 20 turns to find the code.} \\
\\ 
\texttt{What's my first guess? Or shall I start with something like '{[}1 2 3 4{]}'?} \\
            }
        }
    }
    & & \\ \\

    \theutterance \stepcounter{utterance}  
    & & & \multicolumn{4}{p{0.6\linewidth}}{
        \cellcolor[rgb]{0.9,0.9,0.9}{
            \makecell[{{p{\linewidth}}}]{
                \texttt{\tiny{[P1$\langle$GM]}}
                \texttt{{[}GAME{]} Submitted {[}1 2 3 4{]}. Feedback: 1 black peg(s), 2 white peg(s).} \\
            }
        }
    }
     \\ \\

    \theutterance \stepcounter{utterance}  
    & \multicolumn{4}{p{0.6\linewidth}}{
        \cellcolor[rgb]{0.561,0.737,0.561}{
            \makecell[{{p{\linewidth}}}]{
                \texttt{\tiny{[P1$\rangle$GM]}}
                \texttt{Great! I have my first feedback: **1 black peg, 2 white pegs**.} \\
\\ 
\texttt{**Analysis:**} \\
\texttt{{-} **Black pegs (1):** One digit is in the correct position.} \\
\texttt{{-} **White pegs (2):** Two digits are in the correct set of values but in the wrong positions.} \\
\texttt{{-} Since there are 3 digits involved in the feedback (1 correct position + 2 correct values elsewhere), and no duplicates are allowed in the code, the fourth digit I guessed must be wrong (not in the code at all).} \\
\\ 
\texttt{**My Guess:** `{[}1 2 3 4{]}`} \\
\texttt{{-} Position 1 (`1`) is correct (Black).} \\
\texttt{{-} Digits `2` and `4` are in the code but in the wrong slots (White).} \\
\texttt{{-} Digit `3` is not in the code at all.} \\
\\ 
\texttt{**Deductions:**} \\
\texttt{1. The code contains `1` at position 1.} \\
\texttt{2. The code contains the values `2` and `4` (in some order), but not at positions 2 or 4.} \\
\texttt{3. The code does **not** contain `3`.} \\
\texttt{4. Since `2` and `4` are present, and `1` is at position 1, I need to place `2` and `4` in the remaining slots: positions 2, 3, and 4.} \\
\texttt{   {-} Position 1 is fixed as `1`.} \\
\texttt{   {-} I need to figure out the arrangement of `2` and `4` in slots 2, 3, and 4. Wait, I only have two values (`2` and `4`) to fill three slots? No, that's impossible if no duplicates are allowed.} \\
\\ 
\texttt{Let me re{-}evaluate.} \\
\texttt{{-} The code has 4 unique digits.} \\
\texttt{{-} My guess `{[}1 2 3 4{]}` yielded 1 Black and 2 Whites.} \\
\texttt{{-} 1 Black means one digit is correct in its position. Let's assume `{[}1{]}` is correct. So Code = `1 \_ \_ \_`.} \\
\texttt{{-} 2 Whites means two other digits are in the code but in the wrong positions.} \\
\texttt{{-} So, out of `\{1, 2, 3, 4\}`, one is correct (the `1` at pos 1), and two are correct values but misplaced. This accounts for 3 of the 4 digits.} \\
\texttt{{-} The remaining digit must be the one that is neither in the correct position nor present in the wrong positions.} \\
\texttt{{-} Therefore, `3` and `4` are the digits that are present in the code but not at positions 2 or 4.} \\
\texttt{{-} Wait, if `2` is white, it means `2` is in the code but not at position 2.} \\
\texttt{{-} If `4` is white, it means `4` is in the code but not at position 4.} \\
\texttt{{-} And `1` is black, so `1` is at position 1.} \\
\texttt{{-} That means the code contains `1`, `2`, and `4`.} \\
\texttt{{-} The 4th digit must be the remaining digit from the set `\{2, 3, 4, 6\}` (since `1` and `6` and `5` are the other options, but `6` and `5` weren't guessed).} \\
\texttt{{-} Actually, the possible digits are `1, 2, 3, 4, 5, 6`.} \\
\texttt{{-} My guess used `1, 2, 3, 4`.} \\
\texttt{{-} Feedback: 1 Black (`1` is correct at pos 1), 2 Whites (`2` and `4` are in the code but not at pos 2 and pos 4 respectively).} \\
\texttt{{-} This implies the code contains `1`, `2`, and `4`.} \\
\texttt{{-} The 4th digit is not `1`, `2`, or `4`.} \\
\texttt{{-} The remaining available digits are `3`? No, I guessed `3` and it didn't yield a peg. So `3` is not in the code.} \\
\texttt{{-} The remaining available digits are `5` and `6`.} \\
\texttt{{-} Since the code has 4 unique digits, and I have identified `1`, `2`, `4`, and I know `3` is not in the code (because `3` gave no pegs), the 4th digit must be either `5` or `6`.} \\
\texttt{{-} So the code is `1 \_ \_ \_`, with digits `\{1, 2, 4, 5\}` or `\{1, 2, 4, 6\}`.} \\
\texttt{{-} I know `2` is not at position 2.} \\
\texttt{{-} I know `4` is not at position 4.} \\
\texttt{{-} So `2` must be at pos 3 or 4.} \\
\texttt{{-} And `} \\
            }
        }
    }
    & & \\ \\

    \theutterance \stepcounter{utterance}  
    & & & \multicolumn{4}{p{0.6\linewidth}}{
        \cellcolor[rgb]{0.9,0.9,0.9}{
            \makecell[{{p{\linewidth}}}]{
                \texttt{\tiny{[P1$\langle$GM]}}
                \texttt{{[}GAME{]} You attempted an invalid move. Reason: The guess should contain exactly 4 numbers. Please resubmit a valid move and remember to follow the game rules to avoid penalties.} \\
            }
        }
    }
     \\ \\

    \theutterance \stepcounter{utterance}  
    & \multicolumn{4}{p{0.6\linewidth}}{
        \cellcolor[rgb]{0.561,0.737,0.561}{
            \makecell[{{p{\linewidth}}}]{
                \texttt{\tiny{[P1$\rangle$GM]}}
                \texttt{I apologize for the confusion in my previous internal monologue. Let me restart the logic cleanly and submit a valid move immediately.} \\
\\ 
\texttt{**Recap of State:**} \\
\texttt{{-} **Guess:** `{[}1 2 3 4{]}`} \\
\texttt{{-} **Feedback:** 1 Black, 2 White.} \\
\texttt{{-} **Deductions:**} \\
\texttt{  {-} **Black (1):** One digit is correct and in the right position.} \\
\texttt{  {-} **White (2):** Two digits are in the code but in the wrong positions.} \\
\texttt{  {-} **Total unique digits identified:** 1 (Black) + 2 (White) = 3 distinct digits present in the code.} \\
\texttt{  {-} **Digits NOT in the code:** The remaining 1 digit from my guess is not in the code at all.} \\
\\ 
\texttt{**Analysis of `{[}1 2 3 4{]}`:**} \\
\texttt{{-} If I assume `1` is the Black peg (Position 1 is correct):} \\
\texttt{  {-} Then `2` and `4` are the White pegs. This means `2` is in the code (not at Pos 2) and `4` is in the code (not at Pos 4).} \\
\texttt{  {-} `3` is not in the code.} \\
\texttt{  {-} The code contains `\{1, 2, 4, ?\}`. The missing digit must be `5` or `6`.} \\
\texttt{{-} If I assume `2` is the Black peg (Position 2 is correct):} \\
\texttt{  {-} Then `1` and `4` (or `1` and `3`, etc.) are White.} \\
\texttt{  {-} This is also a valid combination of 3 digits.} \\
\\ 
\texttt{To make the most progress quickly, I will test a guess that keeps the **Black peg** (`1` at Pos 1) to verify it, while trying to resolve the positions of the **White pegs** (`2` and `4`).} \\
\\ 
\texttt{**Strategy:**} \\
\texttt{1. Keep `1` at Position 1 (to maintain the Black peg).} \\
\texttt{2. Keep `2` and `4` (since they are confirmed present).} \\
\texttt{3. Move `3` out of the equation (since it yielded no pegs).} \\
\texttt{4. Insert the remaining unknown digit (`5` or `6`) into the slot vacated by `3`.} \\
\texttt{   {-} Since `2` is White, it cannot be at Position 2. It must be at Pos 3 or Pos 4.} \\
\texttt{   {-} Since `4` is White, it cannot be at Position 4. It must be at Pos 2 or Pos 3.} \\
\\ 
\texttt{Let's try **Guess 2: `{[}1 4 2 5{]}`**.} \\
\texttt{{-} Digits: `1, 4, 2, 5`.} \\
\texttt{{-} If `1` is still Black (1 Black).} \\
\texttt{{-} `4` is now at Pos 2. Since `4` was White in the last guess (meaning it was in the code but not at Pos 4), moving it to Pos 2 might turn it into a Black peg.} \\
\texttt{{-} `2` is now at Pos 3. Since `2` was White in the last guess (meaning it was in the code but not at Pos 2), moving it to Pos 3 might turn it into a Black peg.} \\
\texttt{{-} `5` is new.} \\
\\ 
\texttt{This setup is ideal: if both `4` and `2` are correctly placed in their new spots, I should get **3 Black Pegs**. If only one is correct, I get **2 Black, 1 White**. This will drastically narrow down the possibilities.} \\
\\ 
\texttt{**Valid Move:** `{[}1 4 2 5{]}`} \\
\\ 
\texttt{{[}1 4 2 5{]}} \\
            }
        }
    }
    & & \\ \\

    \theutterance \stepcounter{utterance}  
    & & & \multicolumn{4}{p{0.6\linewidth}}{
        \cellcolor[rgb]{0.9,0.9,0.9}{
            \makecell[{{p{\linewidth}}}]{
                \texttt{\tiny{[P1$\langle$GM]}}
                \texttt{{[}GAME{]} Submitted {[}1 4 2 5{]}. Feedback: 1 black peg(s), 1 white peg(s).} \\
            }
        }
    }
     \\ \\

    \theutterance \stepcounter{utterance}  
    & \multicolumn{4}{p{0.6\linewidth}}{
        \cellcolor[rgb]{0.561,0.737,0.561}{
            \makecell[{{p{\linewidth}}}]{
                \texttt{\tiny{[P1$\rangle$GM]}}
                \texttt{**Analysis of Feedback:**} \\
\texttt{{-} **Previous Guess:** `{[}1 2 3 4{]}` → 1 Black, 2 White.} \\
\texttt{  {-} Implies: 3 digits are correct (1 in place, 2 misplaced, 4 misplaced). 1 digit incorrect (`3`).} \\
\texttt{  {-} Current Candidate Set: `\{1, 2, 4\}` plus one unknown from `\{5, 6\}`.} \\
\texttt{{-} **Current Guess:** `{[}1 4 2 5{]}` → 1 Black, 1 White.} \\
\texttt{  {-} The Black peg remains at `1` (Position 1). This confirms `1` is fixed at Pos 1.} \\
\texttt{  {-} The White peg count dropped from 2 to 1.} \\
\texttt{  {-} In the previous guess, `2` (at Pos 2) and `4` (at Pos 4) were both White.} \\
\texttt{  {-} In this guess, `4` is at Pos 2 and `2` is at Pos 3.} \\
\texttt{  {-} Since we only have **1 White peg** now, it means only **one** of the digits `4` or `2` is currently in the correct position. The other one is still misplaced (or we swapped them and only one landed correctly).} \\
\texttt{  {-} Consequently, the new digit `5` is **not** in the code (otherwise we would have had at least 2 Black pegs if both `2` and `4` were correct, or 2 White pegs if neither is correct but `5` is present). Wait, let's refine:} \\
\texttt{    {-} We know `1`, `2`, and `4` are in the code.} \\
\texttt{    {-} We kept `5`. If `5` were in the code, we would have 3 correct digits total (1 Black + 2 Whites or more).} \\
\texttt{    {-} Since we have only 1 Black and 1 White:} \\
\texttt{      {-} We know `1` is correct (Black).} \\
\texttt{      {-} We have 1 White. This accounts for one of the remaining correct digits (`2` or `4`) being in the right spot.} \\
\texttt{      {-} Therefore, the other remaining correct digit is in the wrong spot.} \\
\texttt{      {-} And `5` is **not** in the code.} \\
\\ 
\texttt{**Deductions:**} \\
\texttt{1. Code starts with `1`.} \\
\texttt{2. Digits `2` and `4` are in the code.} \\
\texttt{3. Digit `5` is **NOT** in the code.} \\
\texttt{4. Therefore, the 4th digit of the code must be `6`.} \\
\texttt{   {-} Code composition: `1`, `2`, `4`, `6`.} \\
\texttt{5. Positioning:} \\
\texttt{   {-} Pos 1 = `1` (Confirmed Black).} \\
\texttt{   {-} We need to place `2` and `4` in Pos 2, Pos 3, and Pos 4 (with `6` filling the rest).} \\
\texttt{   {-} From previous guess `{[}1 2 3 4{]}`: `2` was at Pos 2 (White), `4` was at Pos 4 (White).} \\
\texttt{   {-} From current guess `{[}1 4 2 5{]}`: `4` is at Pos 2 (contributed to the single White peg?), `2` is at Pos 3 (not the White peg?).} \\
\texttt{   {-} Since `5` is not in the code, it cannot be the source of the White peg.} \\
\texttt{   {-} So the White peg in `{[}1 4 2 5{]}` must come from either `4` (at Pos 2) or `2` (at Pos 3).} \\
\texttt{   {-} If `4` is at Pos 2 and yields a White peg, then `4` is correct but not at Pos 2? No, if `4` is at Pos 2 and it yields a Black peg, we'd have 2 Blacks (`1` and `4`). But we only have 1 Black. So `4` at Pos 2 is **Wrong**.} \\
\texttt{   {-} If `4` at Pos 2 is Wrong, then `2` at Pos 3 must be the source of the White peg. So `2` is in the code at Pos 3.} \\
\texttt{   {-} If `2` is at Pos 3, then `4` must be at Pos 4 (the last slot).} \\
\\ 
\texttt{**Proposed Logic:**} \\
\texttt{{-} Pos 1: `1` (Black)} \\
\texttt{{-} Pos 2: `4` (Misplaced previously, so still misplaced? Or did it move?) {-}> If `4` was White at Pos 4 previously, and now it's at Pos 2, and we still have only 1 White, then `4` is NOT at Pos 2. So `4` must be at Pos 3 or 4.} \\
\texttt{{-} Let's re{-}evaluate the} \\
            }
        }
    }
    & & \\ \\

    \theutterance \stepcounter{utterance}  
    & & & \multicolumn{4}{p{0.6\linewidth}}{
        \cellcolor[rgb]{0.9,0.9,0.9}{
            \makecell[{{p{\linewidth}}}]{
                \texttt{\tiny{[P1$\langle$GM]}}
                \texttt{{[}GAME{]} You attempted an invalid move. Reason: You have already guessed {[}1, 4, 2, 5{]}. Please try a different guess. Please resubmit a valid move and remember to follow the game rules to avoid penalties.} \\
            }
        }
    }
     \\ \\

    \theutterance \stepcounter{utterance}  
    & \multicolumn{4}{p{0.6\linewidth}}{
        \cellcolor[rgb]{0.561,0.737,0.561}{
            \makecell[{{p{\linewidth}}}]{
                \texttt{\tiny{[P1$\rangle$GM]}}
                \texttt{I see. I need to avoid repeating previous guesses. So far:} \\
\texttt{{-} Guess 1: `{[}1 2 3 4{]}`} \\
\texttt{{-} Guess 2: `{[}1 4 2 5{]}`} \\
\\ 
\texttt{I cannot guess `{[}1 4 2 5{]}` again, nor can I guess `{[}1 2 3 4{]}`.} \\
\\ 
\texttt{**Recap of Deductions:**} \\
\texttt{{-} **Pos 1** is `1` (Confirmed Black in both guesses).} \\
\texttt{{-} **Digit 5** is NOT in the code (Eliminated).} \\
\texttt{{-} **Digits 2 and 4** are in the code.} \\
\texttt{{-} **Pos 1** = `1`.} \\
\texttt{{-} **Pos 2** contains `4` in Guess 2 (Feedback: 1 Black, 1 White). Since `1` accounts for the Black, the White must come from `2` or `4`.} \\
\texttt{  {-} If `4` were at Pos 2 in Guess 2 and it was Black, the feedback would be 2 Black (unless `1` wasn't Black, but `1` is fixed).} \\
\texttt{  {-} Since feedback is only 1 Black, `4` at Pos 2 is **incorrect** (misplaced).} \\
\texttt{  {-} Therefore, `4` must be at Pos 3 or Pos 4.} \\
\texttt{{-} In Guess 1 (`{[}1 2 3 4{]}`), `2` was at Pos 2 (White) and `4` was at Pos 4 (White). This confirms `2` and `4` are in the code but not at those specific positions.} \\
\texttt{{-} In Guess 2 (`{[}1 4 2 5{]}`), `4` was at Pos 2 (No Black contribution) and `2` was at Pos 3 (White contribution). This confirms `2` is now at Pos 3.} \\
\texttt{{-} So: **Pos 1 = 1**, **Pos 3 = 2**.} \\
\texttt{{-} The remaining slots are Pos 2 and Pos 4.} \\
\texttt{{-} We know `4` is in the code but NOT at Pos 2 (from Guess 2 feedback) and NOT at Pos 4 (from Guess 1 feedback).} \\
\texttt{{-} Wait, if `4` is not at Pos 2 (Guess 2) and not at Pos 4 (Guess 1), where can `4` be?} \\
\texttt{{-} The remaining slot is **Pos 4**.} \\
\texttt{{-} Let's verify: If Code is `{[}1 4 2 4{]}`, that has duplicates. Invalid.} \\
\texttt{{-} Let's re{-}evaluate the "White" definition. White = Correct digit, wrong position.} \\
\texttt{{-} Guess 1: `{[}1 2 3 4{]}`. Feedback 1B, 2W.} \\
\texttt{  {-} 1B: `1` is correct at Pos 1.} \\
\texttt{  {-} 2W: `2` and `4` are correct values, wrong positions. `3` is not in code.} \\
\texttt{  {-} So `2` is in code (not at Pos 2). `4` is in code (not at Pos 4).} \\
\texttt{{-} Guess 2: `{[}1 4 2 5{]}`. Feedback 1B, 1W.} \\
\texttt{  {-} 1B: `1` is correct at Pos 1. (Confirmed).} \\
\texttt{  {-} 1W: One of the other digits is correct value, wrong position.} \\
\texttt{  {-} We know `5` is not in code.} \\
\texttt{  {-} So the White must come from `4` or `2`.} \\
\texttt{  {-} In Guess 2, `4` is at Pos 2. `2` is at Pos 3.} \\
\texttt{  {-} If `4` is at Pos 2 and yields White, then `4` is in code but NOT at Pos 2. (Matches Guess 1 deduction).} \\
\texttt{  {-} If `2` is at Pos 3 and yields White, then `2` is in code but NOT at Pos 3.} \\
\texttt{  {-} Since we only have **1 White** total:} \\
\texttt{    {-} Either (`4` at Pos 2 is White AND `2` at Pos 3 is NOT White) OR (`4` at Pos 2 is NOT White AND `2` at Pos 3 is White).} \\
\texttt{  {-} From Guess 1: `2` is NOT at Pos 2. `4` is NOT at Pos 4.} \\
\texttt{  {-} Scenario A: `2` is the White peg in Guess 2.} \\
\texttt{    {-} This implies `2` is in code but NOT at Pos 3.} \\
\texttt{    {-} But we know `2` is in code. So `2` must be at Pos 4.} \\
\texttt{    {-} Check consistency with Guess 1: Guess 1 said `2` is not at Pos 2. (Consistent} \\
            }
        }
    }
    & & \\ \\

    \theutterance \stepcounter{utterance}  
    & & & \multicolumn{2}{p{0.3\linewidth}}{
        \cellcolor[rgb]{0.9,0.9,0.9}{
            \makecell[{{p{\linewidth}}}]{
                \texttt{\tiny{[GM$|$GM]}}
                \texttt{role: Player 0} \\
\texttt{invalid\_move: True} \\
\texttt{turn\_count: 3} \\
\texttt{reason: Invalid Move: You have already guessed {[}1, 4, 2, 5{]}. Please try a different guess.} \\
\texttt{reward: 0.375} \\
            }
        }
    }
    & & \\ \\

\end{supertabular}
}

\end{document}
